\section{Introduction}
Quantum chromodynamics (QCD) describes quark-gluon interactions over
a wide energy range, posing nonperturbative challenges at lower energies.
Hadron structure studies depend on parton distribution functions (PDFs),
with PDF determinations relying on precise QCD calculations,
influencing investigations at the LHC and future projects (EIC, HL-LHC, LHeC).
The quest for new physics hinges on accurate experimental-theoretical comparisons,
emphasizing PDF precision.

QCD factorization in collider processes, beyond PDFs,
involves generalized parton distributions (GPDs) and
transverse-momentum-dependent distributions (TMDs).
While not directly measurable, these distribution functions
can be determined by factorizing the process's cross section
into hard and soft components.

In deep inelastic scattering (DIS), a key role is played by the structure
functions, $F_k(x,Q^2)~(k=1,2,3)$, which depend on $Q^2=-q^2$, where $q$ is the 4-momentum
transferred to the nucleon, and $x = Q^2/2(pq)$ represents
the fraction of nucleon momentum carried by the parton.
Here, $p$ denotes the nucleon momentum.
For $Q^2$ much larger than the nucleon mass,
the factorization theorem~\cite{Collins:1989gx}
relates the structure functions to the convolution of
the scale-dependent parton distribution functions $f_a(x,\mu^2)$
with the Wilson coefficients $C_k^a(Q^2/\mu^2)$
(summed over quark flavors and gluons, $a$).
The Wilson coefficients are calculable in perturbation theory, and the
evolution of the PDFs in $\mu$ is governed by the
Dokshitzer-Gribov-Lipatov-Altarelli-Parisi (DGLAP)
equations~\cite{Gribov:1972ri,Lipatov:1974qm,Dokshitzer:1977sg,Altarelli:1977zs}.
The DGLAP equations can be solved once we specify an initial condition at
a certain scale $\mu=\mu_0$. The evolution of the PDFs is calculable
in perturbation theory, but the precision of experimental data
always necessitates higher-order calculations.
The standard choice for renormalization
and factorization scheme is the $\MSbar$ scheme.
PDFs can be determined experimentally by analyzing a set of
hard-scattering processes involving nucleons.
There is a long-standing community effort to determine
PDFs using well-defined parametrizations for their $x$-dependence
and starting the evolution around $\mu_0^2 = 1 - 4$GeV.

Lattice QCD can offer a theoretical input for the determination of the
PDFs and their evolution.
The direct calculation of PDFs using lattice QCD poses particular challenges
due to the Euclidean geometry and the light-cone dominance of the kinematics.
The connection between PDFs and
hadronic matrix elements, which are calculable in lattice QCD, is
established through the
moments of the PDFs, denoted as
$\left\langle x^{n} \right\rangle$~\cite{Curci:1980uw,Collins:1981uw}.
Lattice QCD calculations of the moments of the PDFs,
pioneered in Refs.~\cite{Kronfeld:1984zv,Martinelli:1987zd,Martinelli:1987bh},
provide, in principle, a means for the
complete reconstruction of the PDFs.
This possibility has remained impractical
due to the theoretical and numerical challenges
associated with computing high moments.
Alternative approaches have been developed to determine
the $x$-dependence of the PDFs, the heavy-quarks OPE~\cite{Aglietti:1998mz,Detmold:2005gg},
the quasi-PDF~\cite{Ji:2013dva},
the psuedo-PDF~\cite{Radyushkin:2016hsy,Karpie:2018zaz},
the OPE-based method of~\cite{Chambers:2017dov},
the current-current approach~\cite{Braun:2007wv,Ma:2014jla},
and the hadron tensor method of~\cite{Liang:2019frk}.
It is worth noticing that these approaches allow
in principle an indirect determination of the
moments of PDF of nucleons
~\cite{HadStruc:2021qdf,Gao:2022uhg} and pions~\cite{Gao:2020ito,Gao:2022iex}.

Different high-energy scatterings are associated with different PDFs:
unpolarized, helicity, and transversity.
Certain kinematical regions, as well as specific PDFs,
pose challenges for extraction from experimental data.
This emphasizes the significance
of reconstructing the PDFs from lattice QCD simulation data.

In this work, we revisit the possibility of calculating moments of PDFs,
which, if successful, would provide an important and complementary
approach to determine distribution functions from QCD.
We describe a method that addresses both the theoretical and numerical
challenges faced in the past, which hindered the calculation of
moments of any order from lattice QCD.
